% Generated draft: requires mathematical and source review.
\documentclass[11pt]{article}
\usepackage[margin=25mm]{geometry}
\usepackage{fontspec}
\setmainfont{Latin Modern Roman}
\usepackage{amsmath,amssymb,mathtools,unicode-math}
\setmathfont{Latin Modern Math}
\usepackage{enumitem,needspace,xurl,hyperref}
\setlength{\parindent}{0pt}
\setlength{\parskip}{5pt}
\setlength{\emergencystretch}{3em}
\newcommand{\N}{\mathbb{N}}
\newcommand{\Z}{\mathbb{Z}}
\newcommand{\Q}{\mathbb{Q}}
\newcommand{\R}{\mathbb{R}}
\newcommand{\C}{\mathbb{C}}
\newcommand{\F}{\mathbb{F}}
\newcommand{\A}{\mathbb{A}}
\newcommand{\PP}{\mathbb{P}}
\newcommand{\E}{\mathbb{E}}
\newcommand{\eps}{\varepsilon}
\newcommand{\dd}{\,\mathrm{d}}
\begin{document}
\section*{UnsolvedMath ID: 20000808}
\textbf{Source problem code:} catalogue entry
\subsection{Definitions and mathematical statement}
For an algebraically closed field \(k\) of characteristic zero, let \(\operatorname{Hilb}^{D}(\A_k^n)\) denote the Hilbert scheme of length-\(D\) zero-dimensional subschemes of affine \(n\)-space. As specified in the supplied source verification, ``component'' means an irreducible component. A component is called rational when its reduced underlying integral variety is birational over \(k\) to affine space, equivalently its function field is purely transcendental over \(k\). The original question is the existential assertion
\[
 \exists D,n\ge1\ \ \exists W\text{ an irreducible component of }
 \operatorname{Hilb}^{D}(\A_k^n),\qquad W_{\mathrm{red}}\text{ is not rational}.
\]
The original extracted text does not itself specify a ground field; characteristic zero and algebraic closedness are the setting of the supplied answer. The trailing text ``\(>12\) RECORDED ON WHITEBOARD BY IZZET COSKUN, AND EDITED BY LI LI'' is an OCR/footer artifact, not an inequality restricting \(D\) or \(n\).

The supplied aid reports a positive answer: Farkas--Pandharipande--Sammartano constructed components that are not rationally connected for ambient dimensions \(n\ge12\), and Wu's June 2026 one-variable-frame preprint extends this to every \(n\ge10\). No assertion of optimality of these thresholds is supplied.

Its proposed strengthening has the quantifier order
\[
 \forall g\ge24\ \ \exists D_g\ge1\ \ \forall n\ge10\ \ \exists W_{g,n}
 \subset\operatorname{Hilb}^{D_g}(\A_k^n),
\]
where \(W_{g,n}\) is an irreducible component whose reduced variety admits a dominant rational map to the moduli space \(\mathcal M_g\) of smooth projective connected genus-\(g\) curves, and whose maximal rationally connected quotient has dimension at least \(3g-3\). Here MRC dimension means the dimension of the base of the maximal rationally connected fibration of a smooth projective model of the function field. The aid derives the MRC conclusion using functoriality and the properties of the genus-\(g\) moduli target; domination alone without a suitable non-uniruled target would not establish that bound. It claims the length can be kept unchanged while extending ambient dimension, so the same \(D_g\) works for every \(n\ge10\).

The reported consequences are that these reduced components are neither rational, stably rational, nor unirational, and that MRC dimensions are unbounded already for \(n=10\) and \(n=11\). Stable rationality means becoming rational after taking a product with some affine space; unirationality means admitting a dominant rational map from a projective space. These stronger conclusions and the uniform-length construction are attributed to the unreviewed aid and require verification of the cited results and field hypotheses; no explicit value of \(D_g\) is supplied.
\subsection{Short English statement}
The question asks whether a Hilbert scheme of points has an irreducible component that is not rational. The supplied aid reports such components in characteristic zero, with a June 2026 preprint lowering the ambient-dimension threshold to ten and a stronger uniform-length claim involving moduli of curves.
\subsection{Sources}
\begin{itemize}
\item Local UnsolvedMath registry, permanent ID \(20000808\), source code \texttt{catalogue entry}; \texttt{aim-arithmetic-geometry-notes.json}, \texttt{notes[53]}, tag \texttt{problem}.
\item AIM, \emph{Components of Hilbert Schemes}, four-page problem list, Problem 7, page 1; recorded on whiteboard by Izzet Coskun and edited by Li Li following the workshop of 19--23 July 2010: \url{https://aimath.org/WWN/hilbertschemes/hilbertschemes.pdf}. The supplied verification labels the statement exact and explains the notation. No PDF was accessed for this draft.
\item Farkas--Pandharipande--Sammartano, construction of non-rationally-connected components for \(n\ge12\), attributed by the aid without a paper title, theorem number, or URL.
\item Wu, June 2026 one-variable-frame preprint, attributed threshold \(n\ge10\); no full title, theorem number, or URL is supplied.
\item Research attempt 2: \texttt{full\_solution}, status \texttt{solved}, classification \texttt{SOLVED-IN-LITERATURE}, rigor label \texttt{proved}, novelty confidence low. Registry created and updated 14 August 2026. This machine-generated research aid is not peer reviewed; independently verify before citation.
\item Ulam AI / UnsolvedMath Contributors; CC BY 4.0.
\end{itemize}
\end{document}
